\section{Conclusion}
\label{sec:conclusion}

The question was whether the split-selection correction of conditional
inference improves a feature ranking used for downstream prediction, and what
it costs. On 31 real datasets the answer to the first part is that the correction
changes what the ranking is free of, not how well it predicts: CIF is the best
conditional inference ranker in the study, ahead of cforest, ctree, and single
CIT, and it ranks 3rd of 17 and 2nd of 18 overall, level with random forest,
boosting, and Boruta rankings, behind recursive feature elimination and, on the
high-dimensional panel, behind boosting. The correction removes the cardinality preference of
CART-style splitting, which the NHANES application of the companion article
shows on real clinical predictors, but that removal does not translate into
better top-$k$ prediction on these benchmarks, and out-of-bag permutation
importance removes the same preference from an ordinary random forest.

The answer to the second part is precise. The cost of conditional inference is
the per-threshold Stage~B test, whose permutation budget must grow as
$K^2/\alpha$ in the number of candidate thresholds for the rule to reject at
all; the feature test and its adaptive stopping are not the cost at the
benchmark's budget. The max-type threshold test resolves the same level at
$999K$ evaluations, is calibrated closer to the nominal level, shows no
detectable difference in power at matched size, and moves the downstream
rankings by at most two positions, so it is the configuration a user who wants
conditional inference at a practical cost should choose. For a practitioner
choosing a ranking method for downstream prediction, a random forest or
boosting importance is as good at a small fraction of the cost; conditional
inference earns its cost where the ranking itself must be free of cardinality
bias, and then the threshold test is where the budget goes.
